\section{Diskussion}
\label{sec:Discussion}
Die erste Messung in diesem Experiment ist die Spektrometer-Messung. Sie verdeutlicht, die Sensibilität der äußeren Einwirkungen für das Spektrometer. Mit dem Raumlicht wird, wie in der \autoref{fig:spectra} dargestellt, eine viel höhere Intensität und breiteres Spektrum gemessen. Der Peak vor der Kurve entspricht dem Spektrum von kalt-weißen LEDs. Die Messung ohne Raumlicht zeigt dagegen das Spektrum der Diode etwas genauer, welches im Bereich vom blauen Licht bei etwa $\SI{460}{\nano\meter}$ liegt. Nach dem Abzug der Dark Counts in der \autoref{fig:spectra_dc}, ist eine Verzerrung bei der Messung mit Raumlicht in dem Bereich der äußeren Lichtquelle erkennbar. Dies bestätigt die Voraussetzung die äußerern Einflüsse so gering wie möglich zu halten.\\\\
Zur Messung der radialen Symmetrie zeigt sich, dass die Wahl von den horizontalen Winkeln $h$ von $\SI{-18}{\degree}$ bis $\SI{34.5}{\degree}$ und den vertikalen Winkeln $v$ von $\SI{-6}{\degree}$ bis $\SI{34.5}{\degree}$ als ausreichend erweist, um die radiale Symmetrie in der \autoref{fig:radial} zu erkennen. Dies entspricht der Erwartung bei einer Faser mit einem runden Querschnitt.\\\\
Mithilfe der radialen Symmetrie wird der allgemeine Winkel zur Faserachse $\Theta$ mit der \autoref{eqn:theta} eingeführt, der in der Simulation benutzt wird, um die Intensität der ''Core''- und ''Cladding''-Photonen gegen den Winkel $\Theta$ in der \autoref{fig:SimData} darzustellen. Die Intensität der ''Core''-Photonen steigt dabei zunächst linear an, was durch den Raumwinkel $\text{d}\Omega = \sin{\Theta}\, \text{d}\Theta \text{d}\Phi$ erklärt werden kann. Diese Intensität kommt an ihr Maximum bei dem kritischen Winkel $\Theta_{\text{krit, core}}$, der mit der \autoref{eqn:theta_krit} bestimmt wurde. Danach werden die ''Core''-Photonen unterdrückt, da diese vermehrt in die Ummantelung übergehen. Es werden jedoch noch weitere ''Core''-Photonen über diesen Winkel gemessen, da die Photonen mit größeren Abstand zur Faserachse sich mit mehreren Totalreflexionen an den Grenzflächen durch die Faser bewegen und der tatsächliche Reflexionswinkel damit kleiner ist, als der Winkel zur Faserachse.\\
Die daraus folgenden ''Cladding''-Photonen werden dann vermehrt detektiert bis diese ebenfalls an einen kritischen Winkel $\Theta_{\text{krit, clad}}$ an ein Maximum der Intensität gelangen. Der weitere Verlauf wird dann stärker unterdrückt, da immer wenigere Laufbahnen die Totalreflexion-Bedinung erfüllen und dementsprechend wenigere Photonen gemessen werden. Der größte Winkel, wo noch Photonen simuliert detektiert werden, liegt bei etwa $\SI{45}{\degree}$, was vermutlich an den Möglichkeiten der Laufbahnen liegt. Der kritische Winkel zum Übergang von der äußeren Ummantelung zur Luftschicht liegt dabei bei etwa $\Theta_{\text{krit, clad2}} \approx \SI{51.3}{\degree}$, wo auch keine Ergebnisse mehr vorliegen.\\
Mithilfe des minimalen Abstandes zur Faserachse mit der \autoref{eqn:r_min} und \autoref{eqn:r_min_parallel} werden diese Ergebnisse in ein zwei-dimensionales Histogramm mit der zusätzlichen Abhängigkeit zu $r_{\text{min}}$ eingezeichnet in der \autoref{fig:hist2d}. Die kritischen Reflexionswinkel $\Theta_{\text{refl}}$ werden mit der \autoref{eqn:skript} ebenfalls visualisiert. Diese zeigen die klare Grenze der gemessenen Photonen und dass diese mit größeren minimalen Abstand zur Faserachse auch größere Winkel für die Totalreflexion erlauben. Zudem nimmt die ''Core''-Intensität auch schon ab dem Bereich für kleine $\Theta$ als auch für kleine $r_{\text{min}}$ zu, währenddessen die ''Cladding''-Intensität erst mit der Grenzfläche von der ''Core''-Intensität gemessen wird. Der Verlauf der ''Cladding''-Photonen hat dabei einen kleinen optischen Schnitt, welcher vermutlich die Grenze an Totalreflexionen an dem Faserende symbolisieren soll.\\
Im Letzten Schritt der Simulation sollen die Dämpfungslängen bestimmt werden. Dies wurde an die Einstellungen der Intensitätsmessung aus dem \autoref{sec:intensity} angepasst, sodass die Winkelabhängigkeit zu $\Theta$ im Bereich von $\SI{0}{\degree} \pm \SI{2}{\degree}$ bis $\SI{36}{{\degree}} \pm \SI{2}{\degree}$ dargestellt wurde. Dabei lieferte die \autoref{fig:atten} die Dämpfungslängen von $\SI{4868}{\milli\meter}$ bis zu $\SI{5328}{\milli\meter}$. Zudem ist erkennbar, dass die effektiven Dämpfungslängen stark varrieren. Das liegt daran, dass sowohl die ''Core''- als auch die ''Cladding''-Photonen unter einer effektiven Größe beschrieben werden. Die Dämpfungslänge liegt im Bereich von wenigen Metern, weshalb die Berechnung realistisch erscheint.\\\\
Die letzte Messung ist die Intensitätmessung, welche die gemessenen Werte aus dem vorherigen Simulationsabschnitt liefert. Die Ergebnisse sind dabei in der \autoref{fig:intensity} dargestellt. Diese liegen mit $\SI{1530}{\milli\meter}$ bis zu $\SI{2182}{\milli\meter}$ unter den simulierten Werten. Klar erkennbar sind dabei zwei deutliche Knicke in dem Verlauf der Kurven, was durch Beschädigungen an der Faser als auch andere äußere Einwirkungen begründet werden könnte. Da diese Werte ebenfalls in dem erwarteten Bereich für die Dämpfungslänge liegen, wird davon ausgegangen, dass ein systematischer Fehler im Bezug auf die Messung auszuschließen ist und durch die obenen genannten Begründungen erklärt werden könnte. Ein Beispiel liefert dabei die Berechnung der wenigen Daten nach dem letzten Knick, die die Dämpfungslängen von $\SI{2261}{\milli\meter}$ bis zu $\SI{3715}{\milli\meter}$ liefern. Dies zeigt, dass wenige Einwirkungen reichen um die Dämpfungslängen stark zu verändern.\\\\
Zusammenfassend konnte dieses Experiment das Prinzip der szintillierenden Fasern in dem LHCb Detektor ansatzweise beschreiben. In diesem Protokoll konnten die Verhaltensweisen der Photonen in einer Faser beschrieben werden und abschließend realistische Werte der Dämpfungslängen der Photonen in einer Faser ermittelt werden.